\documentclass[aspectratio=169]{beamer}

\usetheme{Madrid}
\usecolortheme{default}
\setbeamertemplate{navigation symbols}{}

\usepackage[T1]{fontenc}
\usepackage{lmodern}
\usepackage{booktabs}
\usepackage{array}
\usepackage{graphicx}

\title[Campus Management Database]{Database Design for a Campus Space Management System}
\subtitle{An overview of our design and implementation}
\author{Group 06}
\date{\today}

\begin{document}

\begin{frame}
    \titlepage
\end{frame}

\begin{frame}{What Problem Are We Solving?}
    The system manages the physical spaces of the School of Computer Science.
    Its main goal is to make room usage traceable from the first booking request
    to the final reservation, check-in, or maintenance record.

    \vspace{0.4cm}
    In practical terms, the database must answer questions such as:
    \begin{itemize}
        \item Which spaces are available for a given time period?
        \item Who requested a space, and who approved or rejected it?
        \item What facilities exist in each room?
        \item Which rooms are currently in use or under maintenance?
    \end{itemize}
\end{frame}

\begin{frame}{Main actors of the Database}
    The design separates users by how they interact with the system.

    \begin{itemize}
        \item Users: 
        \begin{itemize} 
            \item Casual End Users
            \item Naive End Users
            \item Sophisticated End Users
        \end{itemize}
        \item Facility staff and facility managers review requests, monitor check-ins, and handle space operations.
        \item Maintenance personnel report, receive, and complete maintenance tasks.
        \item Developers and administrators maintain the database structure, permissions, and reliability.
    \end{itemize}
    \vspace{0.3cm}
    This matters because the database does more than store data. It also enforces
    which people are allowed to perform sensitive actions, especially approvals.
\end{frame}

\begin{frame}{Miniworld Summary}
    The database is centered around four operational scenarios.

    \begin{enumerate}
        \item A user requests a space for a purpose, time slot, and expected number of participants.
        \item A privileged reviewer approves, rejects, cancels, or leaves the request pending.
        \item An approved request becomes a reservation and may later be checked in.
        \item A space may contain facilities and may also be temporarily unavailable because of maintenance.
    \end{enumerate}

    \vspace{0.3cm}
    The design therefore treats spaces, users, booking requests, reservations,
    facilities, policies, and maintenance records as the core operational objects.
\end{frame}

\begin{frame}{Core Operational Entities}
    \begin{table}
        \small
        \begin{tabular}{p{0.24\linewidth} p{0.64\linewidth}}
            \toprule
            \textbf{Entity} & \textbf{Plain-language role in the system} \\
            \midrule
            \texttt{User} & A person who can request, review, attend, or maintain a space. \\
            \texttt{Space} & A physical room or area that can be booked, used, closed, or maintained. \\
            \texttt{Facility} & Equipment located in a space, such as chairs, projectors, boards, or air conditioners. \\
            \texttt{BookingRequest} & A request to use a space during a requested time slot. \\
            \texttt{Review} & The decision record for a request: pending, approved, rejected, or cancelled. \\
            \texttt{Reservation} & The confirmed record created after a request is approved. \\
            \texttt{Maintenance} & A record of a problem, its assigned repair work, and its final result. \\
            \bottomrule
        \end{tabular}
    \end{table}
\end{frame}

\begin{frame}{Reference Entities}
    Many attributes use lookup tables instead of free text.
    This keeps the data consistent and easier to query.

    \vspace{0.3cm}
    Important reference entities include:
    \begin{itemize}
        \item \texttt{SpaceType}, \texttt{SpaceStatus}, and \texttt{SpaceCondition};
        \item \texttt{UserRole}, \texttt{UserStatus}, and \texttt{Department};
        \item \texttt{FacilityType} and \texttt{BookingPurpose};
        \item \texttt{ReviewDecision}, \texttt{ReservationStatus}, and \texttt{MaintenanceStatus}.
    \end{itemize}

    \vspace{0.3cm}
    For example, a \texttt{ReservationStatus} is not typed manually as a sentence.
    It references one controlled value such as \texttt{PENDING}, \texttt{CHECK_IN},
    \texttt{COMPLETED}, or \text{NO_SHOW}.
\end{frame}

\begin{frame}{Conceptual ER Diagram}
    \begin{figure}
            \centering
            \includegraphics[width=0.88\linewidth]{assets/conceptual.jpg}
\end{figure}
\end{frame}

\begin{frame}{Booking and Reservation Flow}
    The booking flow preserves every important step instead of overwriting old data.

    \begin{enumerate}
        \item A requester creates a \texttt{BookingRequest}.
        \item The \texttt{Booking} table connects that request to exactly one user and one space.
        % \item A \texttt{Review} stores the reviewer, decision, decision time, notes, and rejection reason when relevant.
        \item If the request is approved, a \texttt{Reservation} is created from it.
        \item During actual use, \texttt{ReservationCheckin} records the attendant, check-in user, actual times, and room condition.
    \end{enumerate}

    \vspace{0.3cm}
    This structure supports both planning questions and audit questions.
\end{frame}

\begin{frame}{Facilities and Maintenance}
    The space model is also designed for day-to-day operations.

    \begin{itemize}
        \item Each space follows a \texttt{SpacePolicy}, which stores booking window, minimum duration, maximum duration, and check-in grace time.
        \item Each space contains facilities, and each facility has a controlled facility type.
        \item Maintenance begins with a reporter and a description of the problem.
        \item The \texttt{Maintaining} table assigns the maintenance work to a technician, a space, and a maintenance time slot.
        \item A completed maintenance record can store a result note explaining the outcome.
    \end{itemize}
\end{frame}

\begin{frame}{Logical Diagram}
    \begin{figure}
            \centering
            \includegraphics[width=0.86769420\linewidth]{assets/logical.jpg}
\end{figure}
\end{frame}

\begin{frame}{Relational Implementation}
    The logical model is implemented as SQL Server tables.

    \begin{itemize}
        \item Lookup tables are grouped under the \texttt{lookup\_table} schema.
        \item Ternary relationships are implemented as junction tables:
            \texttt{Booking}, \texttt{ReservationCheckin}, and \texttt{Maintaining}.
        \item \texttt{Review} is modeled as a relationship-like table because review decisions have their own attributes.
        \item Primary keys identify each record, while foreign keys connect requests, users, spaces, reservations, and statuses.
    \end{itemize}

    \vspace{0.3cm}
    The result is a database that is normalized enough to avoid repeated status
    text, but still direct enough for common operational queries.
\end{frame}

\begin{frame}{Integrity Rules}
    The design uses constraints and triggers to protect business meaning.

    \begin{itemize}
        \item IDs follow expected formats, such as 8-digit user IDs and fixed-length request or reservation IDs.
        \item Time intervals must be valid: end time must be after start time.
        \item Expected participants must be positive and cannot exceed room capacity.
        \item Only active users can book spaces.
        \item Only facility staff and facility managers can approve or reject requests.
        \item Approved bookings cannot overlap for the same space.
        \item A room under ongoing maintenance must have the correct space status.
    \end{itemize}
\end{frame}

\begin{frame}{Sample Data}
    The sample data turns the schema into a realistic campus scenario.

    \begin{itemize}
        \item Space policies define different booking limits for different kinds of rooms.
        \item Spaces include auditoriums, classrooms, lecture halls, meeting rooms, study rooms, library rooms, staff rooms, and laboratories.
        \item Rooms are distributed across buildings and floors, with capacities from small study rooms to large auditoriums.
        \item Facilities include chairs, air conditioners, projectors, desks, computers, microphones, speakers, and other equipment.
    \end{itemize}

    \vspace{0.3cm}
    This sample data is useful because it lets the team test both ordinary cases
    and edge cases such as unavailable rooms or maintenance conflicts.
\end{frame}

\begin{frame}{Query Layer}
    The stored procedures are written around common business questions.

    \begin{itemize}
        \item Booking views: approved requests after a date, user booking history, pending requests, and requests within a time frame.
        \item Space views: available spaces for a time frame, facilities in a space, spaces with enough capacity, and space utilization.
        \item Reservation views: no-show reservations, upcoming approved bookings, and reservations at a timestamp.
        \item Maintenance views: spaces under maintenance and maintenance assigned to a technician.
        \item Management views: rejection counts, frequent bookers, bad users, and frequent no-show users.
    \end{itemize}
\end{frame}

\begin{frame}{Design Refinements}
    The validation step improved the model in three important ways.

    \begin{itemize}
        \item Department membership was clarified: academic users may belong to departments, but facility-side users may not.
        \item Space policy became its own entity instead of a plain text paragraph.
        \item User status and space condition became reference entities, which makes validation and reporting cleaner.
    \end{itemize}

    \vspace{0.3cm}
    These refinements make the model more precise without making the user workflow
    harder to understand.
\end{frame}

\begin{frame}{What the Database Achieves}
    This database provides a complete operational trail for campus space management.

    \begin{itemize}
        \item It records who asked for a space, what they needed, and when they needed it.
        \item It records who reviewed the request and what decision was made.
        \item It connects approved requests to reservations and actual check-ins.
        \item It tracks space availability, facilities, policies, room condition, and maintenance.
        \item It supports both daily operations and summary reporting through stored procedures.
    \end{itemize}
\end{frame}

\begin{frame}{Closing}
    In short, the design turns campus space management into a clear database workflow:

    \vspace{0.4cm}
    \begin{center}
        request \(\rightarrow\) review \(\rightarrow\) reservation \(\rightarrow\) check-in
    \end{center}

    \vspace{0.4cm}
    Around that workflow, the database also keeps track of spaces, policies,
    facilities, statuses, maintenance, and reporting queries.

    \vspace{0.4cm}
    The main strength of the design is that it does not only store the final state
    of a room. It stores the process that produced that state.
\end{frame}

\end{document}
